\begin{lemma}
\label{editorial-source-lemma-make-separably-generated}
Let $K/k$ be a finitely generated field extension.
There exists a diagram
$$
\xymatrix{
K \ar[r] & K' \\
k \ar[u] \ar[r] & k' \ar[u]
}
$$
where $k'/k$, $K'/K$ are finite purely inseparable field
extensions such that $K'/k'$ is a separably generated field extension.
\end{lemma}

\begin{proof}
Choose the original transcendence basis
$x_1,\ldots,x_r$ of $K$ over $k$, and put
$F=k(x_1,\ldots,x_r)$. Its size $r$ is finite
because the original field extension is finitely
generated. If $K=k(z_1,\ldots,z_m)$ is a specified
finite generating family, then every $z_j$ is
algebraic over $F$ and $K=F(z_1,\ldots,z_m)$.
Successively adjoining these original elements
therefore makes $K/F$ finite. In characteristic
zero every irreducible algebraic polynomial has
nonzero derivative and is separable, so this
basis already separates $K/k$ and one can take
$k'=k$, $K'=K$.

Suppose now that the characteristic is $p>0$.
Let $K/K_{sep}/F$ be the original separable-first
tower from Fields, Lemma \ref{fields-lemma-separable-first}.
Thus $K_{sep}/F$ is finite separable and
$K/K_{sep}$ is finite purely inseparable. Retain
$d=[K:K_{sep}]$. If $d=1$, the original basis
is separating and again no extension is needed.
We construct the strict induction step for $d>1$.

Choose an original $\eta\in K\setminus K_{sep}$.
Pure inseparability gives a least $h\geq1$
such that $\eta^{p^h}\in K_{sep}$. Set
$$
\beta=\eta^{p^{h-1}},\qquad
\alpha=\beta^p=\eta^{p^h}.
$$
Then $\beta\notin K_{sep}$ and
$\alpha\in K_{sep}$. The polynomial $T^p-\alpha$
is irreducible over $K_{sep}$, since a $p$-th
root of $\alpha$ there would equal $\beta$
in the ambient field, contradicting this choice;
see Fields, Lemma \ref{fields-lemma-take-pth-root}.
Hence $[K_{sep}(\beta):K_{sep}]=p$.
Keep the original monic minimal polynomial
$$
P(T)=T^n+a_1T^{n-1}+\cdots+a_n\in F[T]
$$
of $\alpha$. It is separable since
$\alpha\in K_{sep}$; in particular $P'(\alpha)\ne0$.

Here is the omitted finite coefficient construction.
For every original $a_i$, choose complete polynomial
numerator and denominator expressions
$$
a_i=\frac{A_i(x_1,\ldots,x_r)}{B_i(x_1,\ldots,x_r)},
\qquad A_i,B_i\in k[X_1,\ldots,X_r],\quad B_i\ne0.
$$
No common factors are canceled and all coefficients
and summands of these chosen polynomials are kept.
Let $\mathcal C\subset k$ be the finite set of
all their coefficients. Work in one fixed
algebraic closure $\Omega$ of the original $K$,
and choose the unique $p$-th roots there of
the elements of $\mathcal C$ and of the $x_j$.
Set
$$
k'=k(c^{1/p}:c\in\mathcal C),\qquad
y_j=x_j^{1/p},\qquad
F'=k'(y_1,\ldots,y_r).
$$
Each adjoined constant satisfies $U^p-c=0$.
Thus $k'/k$ is finite purely inseparable,
every element of $k'$ has $p$-th power in $k$,
and $[k':k]\leq p^{|\mathcal C|}$.
The last assertion about powers follows for
polynomials and fractions from the characteristic
$p$ identity $(a+b)^p=a^p+b^p$.
The elements $y_1,\ldots,y_r$ are algebraically
independent over $k'$: if
$\sum_\nu d_\nu y^\nu=0$ is a nonzero polynomial
relation, raising it to the $p$-th power gives
$\sum_\nu d_\nu^p x^\nu=0$ over $k$.
At least one coefficient remains nonzero because
Frobenius is injective on a field, contradicting
the original algebraic independence of the $x_j$.

Write the full original expansions as
$A_i(X)=\sum_\nu a_{i,\nu}X^\nu$ and
$B_i(X)=\sum_\nu b_{i,\nu}X^\nu$, with finite
support and nonnegative integer multi-indices.
Their coefficient-root polynomials are
$$
\widetilde A_i(Y)=\sum_\nu a_{i,\nu}^{1/p}Y^\nu,
\qquad
\widetilde B_i(Y)=\sum_\nu b_{i,\nu}^{1/p}Y^\nu.
$$
The denominator is nonzero at the algebraically
independent original $y_j$, and the full identities
are
$$
\widetilde A_i(y)^p=A_i(x),\qquad
\widetilde B_i(y)^p=B_i(x),\qquad
\left(\frac{\widetilde A_i(y)}{\widetilde B_i(y)}\right)^p=a_i.
$$
Put $b_i=\widetilde A_i(y)/\widetilde B_i(y)\in F'$.
This retains each original numerator, denominator,
monomial and coefficient and proves the required
$p$-th power property for every $a_i$, including zero.
The field $F'$ contains $F$ by the original
inclusions and the identities $x_j=y_j^p$.
Moreover $(F')^p\subseteq F$ and
$[F':F]\leq p^{|\mathcal C|+r}$.

Let $L=KF'\subseteq\Omega$ be the actual compositum.
It gives the original commutative diagram
$$
\xymatrix{
K\ar[r]&L\\
k(x_1,\ldots,x_r)\ar[u]\ar[r]&
k'(x_1^{1/p},\ldots,x_r^{1/p})\ar[u].
}
$$
Every arrow is the indicated inclusion in $\Omega$.
The extension $L/K$ is finite purely inseparable
and $L^p\subseteq K$, because the only new
generators are the displayed constant roots
and coordinate roots. In particular
$[L:K]\leq p^{|\mathcal C|+r}$.
The original field $L$ is also finite over $F'$,
as $K$ is finite over $F$.

Put $C=K_{sep}F'$. A finite collection of separable
generators of $K_{sep}/F$ remains separable after
extension to $F'$: their minimal polynomials over
the larger base divide their original separable
polynomials. Thus $C/F'$ is finite separable.
The extension $L/C$ is purely inseparable. To
check this on every element, take finite
generators of $K/K_{sep}$, choose a common
$p$-power placing all of them in $K_{sep}$,
and apply that power to any rational expression
in these generators with coefficients in $C$.
All resulting powers lie in $C$.
It follows that the original separable-first
field $L_{sep}$ over $F'$ is exactly $C$.
Indeed all of $C$ is separable over $F'$, while
an element of $L$ separable over $F'$ is also
separable over $C$, and is purely inseparable
over $C$ by what was just proved. Its minimal
polynomial over $C$ is then both separable and
purely inseparable, hence linear. This proves
the identification $L_{sep}=K_{sep}F'$ without
an assumption about linear disjointness.

The chosen $\beta$ already belongs to $L_{sep}$.
To verify this with the original polynomial, put
$$
Q(T)=T^n+b_1T^{n-1}+\cdots+b_n\in F'[T].
$$
Then $Q(\beta)^p=P(\alpha)=0$, so
$Q(\beta)=0$. Keeping all derivative factors gives
$$
\begin{split}
Q'(\beta)&=n\beta^{n-1}
+\sum_{i=1}^{n-1}(n-i)b_i\beta^{n-i-1},\\
Q'(\beta)^p&=n\alpha^{n-1}
+\sum_{i=1}^{n-1}(n-i)a_i\alpha^{n-i-1}
=P'(\alpha)\ne0.
\end{split}
$$
Here each integer coefficient is its element of
the prime field, whose $p$-th power is itself;
terms divisible by $p$ are retained with their
actual zero coefficients. Thus the minimal
polynomial of $\beta$ over $F'$ is separable:
it divides $Q$, and differentiating that
factorization at $\beta$ shows that its
derivative is nonzero. Consequently
$\beta\in L_{sep}$. Equivalently, Fields,
Lemma \ref{fields-lemma-pth-root} applies to
the original $\alpha$ and $P$ inside the separable
extension $L_{sep}/F'$; its root $\alpha'$
satisfies $(\alpha')^p=\alpha=\beta^p$ and
therefore $\alpha'=\beta$ by injectivity of
Frobenius in $\Omega$. This proves the original
root identification, including its coefficient map.

All entries of the original induction tower are
therefore the actual indicated subfields:
$$
\xymatrix{
K\ar[r]&L\\
K_{sep}(\beta)\ar[r]\ar[u]&L_{sep}\ar[u]\\
K_{sep}\ar[r]\ar[u]&L_{sep}\ar@{=}[u]\\
k(x_1,\ldots,x_r)\ar[r]\ar[u]&
k'(x_1^{1/p},\ldots,x_r^{1/p})\ar[u]\\
k\ar[r]\ar[u]&k'\ar[u].
}
$$
The decrease can be bounded exactly. Since
$K_{sep}(\beta)\subseteq L_{sep}$ and
$L=KL_{sep}$, a basis of $K$ over
$K_{sep}(\beta)$ spans $L$ over $L_{sep}$.
For completeness, its span is the image of
$L_{sep}\otimes_{K_{sep}(\beta)}K\to L$.
This image is a finite-dimensional algebra over
$L_{sep}$, is a domain as a subring of $L$,
and contains both original fields. Any nonzero
element has invertible multiplication on this
finite-dimensional vector space, so the image
is itself a field. It is therefore the
compositum $L$, proving the asserted spanning.
The original tower law now gives
$$
[L:L_{sep}]
\leq[K:K_{sep}(\beta)]
=\frac{[K:K_{sep}]}{[K_{sep}(\beta):K_{sep}]}
=\frac d p<d.
$$
This verifies the strict induction step without
replacing it by an unproved degree comparison.
Applying the same construction to $L/k'$ with
the retained basis $y_1,\ldots,y_r$ completes
the induction. Compositions of the finite
purely inseparable extensions remain finite
and purely inseparable: choose common powers
for finite generating families at successive
stages and compose those powers. Thus the
final inclusions $k\subseteq k''$ and
$K\subseteq L'$ give the required original
diagram, with $k''$ and $L'$ used as the
final $k'$ and $K'$ in the statement.

\medskip\noindent
The same proof provides finite quantitative
information. The original purely inseparable
degree is $d=p^e$ for some integer $e\geq0$.
Indeed adjoining a purely inseparable algebraic
element has degree a power of $p$: its irreducible
polynomial has only one root in an algebraic
closure, and repeated removal of zero derivatives
writes it as a separable irreducible polynomial
in $T^{p^a}$; the separable polynomial must be
linear because there is only one root. A finite
generating tower and the degree law give $d=p^e$.
Every nonterminal step divides the bound for
the remaining inseparable degree by at least $p$.
Consequently the construction terminates after
some $s\leq e$ steps.

Write the actual fields at these steps as
$k_h\subseteq K_h$ for $0\leq h\leq s$,
with $k_0=k$ and $K_0=K$, and retain the
separating-coordinate candidates
$x_j^{1/p^h}$ in the fixed $\Omega$.
At step $h$, let $c_h$ be the number of
coefficients in the finite set used for its
chosen complete rational numerator and
denominator expressions. The construction gives
$$
[k_{h+1}:k_h]\leq p^{c_h},\qquad
[K_{h+1}:K_h]\leq p^{c_h+r},\qquad
k_{h+1}^{,p}\subseteq k_h,\qquad
K_{h+1}^{,p}\subseteq K_h.
$$
Iterating these actual tower maps yields
$$
[k_s:k]\leq p^{\sum_{h=0}^{s-1}c_h},\qquad
[K_s:K]\leq p^{\sum_{h=0}^{s-1}(c_h+r)},\qquad
k_s^{,p^s}\subseteq k,\qquad K_s^{,p^s}\subseteq K.
$$
The terminal extension
$K_s/k_s(x_1^{1/p^s},\ldots,x_r^{1/p^s})$
is separable by the termination criterion.
These are bounds for the original finite
construction, with all constants and coordinate
root adjunctions retained. They do not assert
that the coefficient counts $c_h$ depend only
on $d$, nor that the original transcendence
basis separates $K/k$ before these extensions.
When $r=0$ the coordinate list is empty; when
$s=0$ the sums are empty, the bounds are $1$,
and the original fields already give the result.
\end{proof}